\chapter{Introduction}
\label{ch:introduction}

Imagine a student preparing for an exam at the Faculty of Information Technology at CTU in Prague.
She has a question about a specific diagram on slide 14 of the lecture notes uploaded to the faculty
course system. To ask her peers, she must switch to Discord, copy the name of the file from the
course page, paste it into a message, and hope that her classmates can locate the same slide quickly
enough to give a useful answer. The question floats away in the message history by the next morning,
disconnected from the material it was about and invisible to students who enrol in the course the
following semester. This scenario, repeated across every course and every cohort, illustrates the
fundamental problem that StudySpace is designed to address: the fragmentation of academic tools that
forces students and instructors to operate across multiple disconnected platforms.

% ----------------------------------------------------------------
\section{Motivation}
\label{sec:motivation}

Modern higher education relies on a combination of platforms that rarely communicate with one
another. At a typical faculty, course materials reside in a learning management system such as Moodle
or the faculty course portal, while student collaboration happens informally on Discord, WhatsApp, or
Microsoft Teams. Code assignments and grading are handled by yet another system, and AI tools such
as ChatGPT are used in isolation without any link to the actual curriculum content. The consequence
of this fragmentation is that information ends up scattered: a useful correction that a student
discovers during self-study has no path back into the course, a question about a specific document
carries no reference to the document itself, and knowledge generated by one cohort is invisible to
the next.

StudySpace addresses these gaps by combining four capabilities in a single environment. Instructors
need a structured way to organise and publish course materials and to manage which students can
access them (F1). Students, rather than passively consuming that content, need a mechanism to clone
a workspace, enrich it with their own notes, and contribute improvements back to the original course
page for the instructor to review and optionally publish (F2). When students want to ask a question
about a specific document, they need a messaging channel that can attach a direct reference to that
document, so other readers can follow the link and land on the relevant material immediately (F3).
Finally, students need an AI assistant that queries the actual course content rather than operating
on general knowledge alone, grounding answers in the material the instructor has provided (F4). All
of these capabilities require a role-based access model that distinguishes instructors from students
at the API level, preventing students from modifying course materials they do not own (F5).

% ----------------------------------------------------------------
\section{Structure}
\label{sec:structure}

Chapter~\ref{ch:analysis} defines the target audience and their pain points, establishes the
functional and non-functional requirements derived from the motivation above, examines existing
platforms and their limitations, and describes the planned system and logical architecture.
Chapter~\ref{ch:design} translates the analysis into concrete technical decisions, covering the
domain model, the technology choices and the rationale behind each, the database schema, and the
user interface design.
Chapter~\ref{ch:implementation} documents the implementation of the five core features, presents
the code listings that illustrate the most important design choices, and describes how the system
was tested.
Chapter~\ref{ch:conclusion} summarises the outcomes relative to the goals stated in this chapter,
reflects on the known limitations of the prototype, and proposes directions for future work.
